% SPDX-License-Identifier: CC-BY-4.0
\documentclass[11pt,letterpaper]{article}

\usepackage[T1]{fontenc}
\usepackage{amsmath,amssymb}
\usepackage{tikz}
\usepackage{pgfplots}
\pgfplotsset{compat=1.18}
\usepackage{algorithm}
\usepackage{algpseudocode}
\usepackage{booktabs}
\usepackage{enumitem}
\usepackage[hidelinks]{hyperref}

\newcommand{\N}{\mathbb{N}}
\newcommand{\Z}{\mathbb{Z}}
\newcommand{\Q}{\mathbb{Q}}
\newcommand{\R}{\mathbb{R}}
\newcommand{\C}{\mathbb{C}}
\newcommand{\Map}[3]{#1 \colon #2 \to #3}
\newcommand{\Mod}[1]{#1\text{-}\mathrm{Mod}}
\newcommand{\GroupHom}{\mathrm{Hom}_{\mathbf{Grp}}}
\newcommand{\RingHom}{\mathrm{Hom}_{\mathbf{Ring}}}

\newtheorem{theorem}{Theorem}
\newtheorem{lemma}[theorem]{Lemma}
\newtheorem{definition}[theorem]{Definition}
\newtheorem{remark}[theorem]{Remark}
\newenvironment{proof}{\noindent\textit{Proof.}}{\hfill\(\square\)}

\title{Coprime scaling}
\author{ascendnoosphere}
\date{\today}

\begin{document}
\maketitle

\section{Coprime scaling and offset normalization}

% All state and parameter signs remain part of the declared dynamics.
\begin{definition}
For integers \(q,c,d\) with \(|d|\ge2\), write
\[
F_{q,c,d}(n)=
\begin{cases}
n/d,&d\mid n,\\
qn+c,&d\nmid n.
\end{cases}
\]
Its state space is \(\Z\). The division branch takes precedence at zero.
For a scale \(h\in\Z\), scaling the parameters means replacing \(c\) by \(hc\),
without changing \(q\) or the signed divisor \(d\).
\end{definition}

% Coprimality is used on the branch guard before comparing arithmetic values.
\begin{lemma}[QCD.Scaling.001.1a]
If \(\gcd(|h|,|d|)=1\), then, for every \(n\in\Z\),
\[
d\mid hn\quad\Longleftrightarrow\quad d\mid n.
\]
\end{lemma}
\begin{proof}
The reverse implication follows by multiplication. For the forward implication,
B\'ezout's identity gives integers \(u,v\) such that \(uh+vd=1\)
\cite{Lang02}. Multiplying by \(n\) gives \(n=u(hn)+v(dn)\).
Both terms on the right are divisible by \(d\), so \(d\mid n\).
\end{proof}

% Both branches and every raw return time are preserved, not just rational values.
\begin{theorem}[QCD.Scaling.001.1]
Suppose \(\gcd(|h|,|d|)=1\). Then
\[
F_{q,hc,d}(hn)=hF_{q,c,d}(n)
\]
for every integer \(n\). Consequently, for every \(s\in\N\),
\[
F_{q,hc,d}^{\,s}(hn)=hF_{q,c,d}^{\,s}(n).
\]
Moreover, \(h\ne0\), and
\[
F_{q,hc,d}^{\,s}(hn)=hn
\quad\Longleftrightarrow\quad
F_{q,c,d}^{\,s}(n)=n.
\]
In particular, the complete sets of return times and primitive raw periods agree.
\end{theorem}
\begin{proof}
The guard lemma makes the two maps choose the same branch. On the division branch,
write \(n=da\); then \((hn)/d=ha=h(n/d)\), including negative \(d\).
On the affine branch,
\[
q(hn)+hc=h(qn+c).
\]
Induction on \(s\), starting at \(s=0\), proves the iterate identity.
If \(h=0\), the coprimality hypothesis would give \(|d|=1\), contradicting
\(|d|\ge2\). Multiplication by the nonzero integer \(h\) is injective, which proves
the return-time equivalence. The least positive return time therefore agrees
whenever such a return time exists.
\end{proof}

% The event count follows the chronological first-step recursion of the actual orbit.
\begin{definition}
For fixed parameters, abbreviate \(F=F_{q,c,d}\). Let \(A(t,n)\) count the affine
instructions among the first \(t\) raw steps from \(n\). Equivalently,
\[
A(0,n)=0,\qquad
A(t+1,n)=
\begin{cases}
A(t,F(n)),&d\mid n,\\
A(t,F(n))+1,&d\nmid n.
\end{cases}
\]
This counts affine events, not raw time or accelerated steps.
\end{definition}

% The division coordinate of the entry rank strictly decreases whenever it is used.
\begin{lemma}[QCD.Scaling.001.2a]
If \(n\ne0\) and \(d\mid n\), then
\[
|F(n)|=|n|/|d|<|n|.
\]
\end{lemma}
\begin{proof}
Write \(n=da\). Since \(n\ne0\), both \(d\) and \(a\) are nonzero, and
\(|n|=|d||a|\). Since \(|d|\ge2\), the quotient \(|a|\) is strictly smaller than
\(|n|\).
\end{proof}

% The prepared divisibility condition survives division and consumes one affine event.
\begin{theorem}[QCD.Scaling.001.2]
Suppose \(\gcd(|c|,|d|)=1\). If \(r\in\N\) and \(c\mid q^r n\), there exists
\(t\in\N\) such that
\[
c\mid F^t(n),\qquad A(t,n)\le r.
\]
\end{theorem}
\begin{proof}
Use well-founded induction on \((r,|n|)\in\N^2\), ordered lexicographically.
If \(c\mid n\), choose \(t=0\). Otherwise \(n\ne0\), because every integer divides
zero. The case \(r=0\) is impossible: its hypothesis would already give
\(c\mid n\). Write \(r=j+1\).

If \(d\mid n\), write \(n=dm\). The hypothesis gives
\(c\mid d q^{j+1}m\). Coprimality permits cancellation of \(d\), so
\(c\mid q^{j+1}m\). The first coordinate of the rank is unchanged, while
\(|m|<|n|\) by the magnitude lemma. Induction gives \(s\) with
\(c\mid F^s(m)\) and \(A(s,m)\le j+1\). Choose \(t=s+1\); the first instruction
is division, so it adds no affine event.

If \(d\nmid n\), put \(m=qn+c\). Then
\[
q^j m=q^{j+1}n+q^j c
\]
is divisible by \(c\). The first coordinate decreases from \(j+1\) to \(j\),
regardless of how large \(|m|\) becomes. Induction gives \(s\) with
\(c\mid F^s(m)\) and \(A(s,m)\le j\). Again choose \(t=s+1\); this time
\(A(t,n)=A(s,m)+1\le j+1\).

In both branches, \(F^{s+1}(n)=F^s(F(n))=F^s(m)\). Each recursive call decreases
the declared rank, and the case \(c\mid n\) supplies the base case.
\end{proof}

% Entry is followed by an exact full-sequence correspondence with offset one.
\begin{theorem}[QCD.Scaling.001.3]
Suppose \(\gcd(|c|,|d|)=1\), \(k\in\N\), and \(c\mid q^k\). For every
\(n\in\Z\), there are \(t\in\N\) and \(m\in\Z\) such that
\[
F_{q,c,d}^{\,t}(n)=cm,\qquad A(t,n)\le k,
\]
and, for every \(s\in\N\),
\[
F_{q,c,d}^{\,t+s}(n)=cF_{q,1,d}^{\,s}(m).
\]
\end{theorem}
\begin{proof}
The power divisibility hypothesis gives \(c\mid q^k n\). Budgeted entry supplies
\(t\) with \(c\mid F_{q,c,d}^{\,t}(n)\) and \(A(t,n)\le k\). Thus write that
state as \(cm\). Apply the scaling theorem to \((q,1,d)\) with scale \(c\), and
split the raw time into the initial \(t\) steps and the following \(s\) steps.
\end{proof}

% The two directions retain every seed and a single common capture witness.
\begin{theorem}[QCD.Scaling.001.4]
Under the hypotheses of the unit-offset tail theorem, let \(T\subseteq\Z\) be any
set. The following obligations are equivalent:
\begin{enumerate}[label=(\roman*)]
\item For every \(m\in\Z\), some raw time \(s\) satisfies
\(cF_{q,1,d}^{\,s}(m)\in T\).
\item For every \(n\in\Z\), some iterate \(F_{q,c,d}^{\,u}(n)\) belongs to \(T\).
\end{enumerate}
No invariance assumption is imposed on \(T\). In particular, for any
\(S\subseteq\Z\), all-seed capture in \(cS\) is equivalent to all-seed
unit-offset capture in \(S\).
\end{theorem}
\begin{proof}
Assume (i). For the tail witnesses \(t,m\) supplied by the preceding theorem,
choose \(s\) with \(cF_{q,1,d}^{\,s}(m)\in T\). At the single raw time \(u=t+s\),
the tail identity gives \(F_{q,c,d}^{\,u}(n)\in T\).

Conversely, assume (ii) and apply it to the particular seed \(cm\). At its capture
time \(u\), the scaling theorem gives
\[
cF_{q,1,d}^{\,u}(m)=F_{q,c,d}^{\,u}(cm)\in T.
\]
This proves (i) without any target invariance or pointwise reversal before entry.

For \(T=cS\), membership \(cx\in cS\) means \(cx=cy\) for some \(y\in S\).
Since \(c\ne0\), cancellation is equivalent to \(x=y\in S\), which proves the
last assertion.
\end{proof}

% Positivity shows that normalization remains inside the original positive domain.
\begin{theorem}[QCD.Scaling.001.5]
Suppose \(q\ge0\), \(c>0\), and \(d>0\). Every positive state has a positive
successor under \(F_{q,c,d}\). Under the additional hypotheses of the unit-offset
tail theorem, a positive seed has a tail witness with \(m>0\). Consequently, the
capture reduction remains valid after restricting both universal quantifiers to
positive seeds.
\end{theorem}
\begin{proof}
If \(d\mid n\), write \(n=da\). Positivity of \(n\) and \(d\) gives \(a>0\), so
the division successor is positive. Otherwise \(qn\ge0\), and hence
\(qn+c>0\). Induction proves positivity of every raw iterate.

For the entry witness \(F_{q,c,d}^{\,t}(n)=cm\), both the left side and \(c\) are
positive, so \(m>0\). The forward capture argument therefore applies its
unit-offset premise only to a positive seed. In the reverse direction, a positive
unit-offset seed \(m\) corresponds to the positive original seed \(cm\). The same
scaling identity proves the target equivalence.
\end{proof}

\section{The reduction and its boundary}

For \(|c|>1\), the capture reduction replaces a nonunit offset by offset one while
retaining \(q\), the signed divisor \(d\), every integer seed, the exact target
correspondence \(T\leftrightarrow\{m:cm\in T\}\), and the raw clock after
entry. Its additional work is the terminating entry calculation. The decreasing
resource in that calculation is the lexicographic rank \((r,|n|)\), not
monotonicity of the raw orbit.
The replacement obligation is capture for \(F_{q,1,d}\); the theorem does not assert
that this replacement obligation is true. For \(|c|=1\), the tail is already
available at time zero, and scaling is an identity or reflection rather than
an offset decrease. The case \(c=0\) is excluded by coprimality and \(|d|\ge2\).

Coprimality is a preservation condition, not a guard that may be discarded.
For example, for \((q,c,d)=(3,1,2)\), scaling by \(h=2\) gives
\[
F_{3,2,2}(2)=1\ne8=2F_{3,1,2}(1).
\]
Thus arbitrary integer scaling is not a raw semiconjugacy.

\bibliographystyle{amsalpha}
\bibliography{Scaling}
\end{document}
